\section{第 8 章综合习题}

\begin{problem}{点在平面上的充要条件}{Point-on-Plane-Condition}
    设 $O$ 为一定点, $A, B, C$ 为不共线的三点. 证明: 点 $M$ 位于平面 $ABC$ 上的充分必要条件是存在实数 $k_1, k_2, k_3$ 使得
    \begin{equation}
        \overrightarrow{OM} = k_1\overrightarrow{OA} + k_2\overrightarrow{OB} + k_3\overrightarrow{OC}, \text{ 且 } k_1 + k_2 + k_3 = 1.
    \end{equation}
\end{problem}

\begin{myproof}
    必要性：若点 $M$ 位于平面 $ABC$ 上，由于 $A, B, C$ 不共线，则向量 $\overrightarrow{AB}$ 与 $\overrightarrow{AC}$ 不共线．
    故存在实数 $x, y$ 使得
    \begin{equation}
        \overrightarrow{AM} = x\overrightarrow{AB} + y\overrightarrow{AC}
    \end{equation}
    依据向量的减法规则
    \begin{equation}
        \overrightarrow{OM} - \overrightarrow{OA} = x(\overrightarrow{OB} - \overrightarrow{OA}) + y(\overrightarrow{OC} - \overrightarrow{OA})
    \end{equation}
    整理得
    \begin{equation}
        \overrightarrow{OM} = (1 - x - y)\overrightarrow{OA} + x\overrightarrow{OB} + y\overrightarrow{OC}
    \end{equation}
    令 $k_1 = 1 - x - y, k_2 = x, k_3 = y$，则有
    \begin{equation}
        \overrightarrow{OM} = k_1\overrightarrow{OA} + k_2\overrightarrow{OB} + k_3\overrightarrow{OC}
    \end{equation}
    且满足
    \begin{equation}
        k_1 + k_2 + k_3 = (1 - x - y) + x + y = 1
    \end{equation}
    充分性：若存在 $k_1, k_2, k_3$ 满足题设条件，由 $k_1 = 1 - k_2 - k_3$ 代入得
    \begin{equation}
        \begin{aligned}
            \overrightarrow{OM}                                & = (1 - k_2 - k_3)\overrightarrow{OA} + k_2\overrightarrow{OB} + k_3\overrightarrow{OC}            \\
            \implies \overrightarrow{OM} - \overrightarrow{OA} & = k_2(\overrightarrow{OB} - \overrightarrow{OA}) + k_3(\overrightarrow{OC} - \overrightarrow{OA}) \\
            \implies \overrightarrow{AM}                       & = k_2\overrightarrow{AB} + k_3\overrightarrow{AC}
        \end{aligned}
    \end{equation}
    由于 $\overrightarrow{AM}$ 可由平面 $ABC$ 内的两个不共线向量 $\overrightarrow{AB}, \overrightarrow{AC}$ 线性表示，故点 $M$ 必在平面 $ABC$ 上．
\end{myproof}


\begin{problem}{向量线性相关性}{Vector-Linear-Dependence}
    证明: 向量 $\symbfit{a} - \symbfit{b} + \symbfit{c}, -2\symbfit{a} + 3\symbfit{b} - 2\symbfit{c}, 2\symbfit{a} - \symbfit{b} + 2\symbfit{c}$ 线性相关.
\end{problem}

\begin{myproof}
    设 $\symbfit{v}_1 = \symbfit{a} - \symbfit{b} + \symbfit{c}$, $\symbfit{v}_2 = -2\symbfit{a} + 3\symbfit{b} - 2\symbfit{c}$, $\symbfit{v}_3 = 2\symbfit{a} - \symbfit{b} + 2\symbfit{c}$．
    考虑其系数矩阵的行列式（以 $\symbfit{a}, \symbfit{b}, \symbfit{c}$ 为基）：
    \begin{equation}
        D = \begin{vmatrix} 1 & -2 & 2 \\ -1 & 3 & -1 \\ 1 & -2 & 2 \end{vmatrix}
    \end{equation}
    观察发现第一列与第三列相同（或第一行与第三行相同），根据行列式性质可知
    \begin{equation}
        D = 0
    \end{equation}
    这表明 $\symbfit{v}_1, \symbfit{v}_2, \symbfit{v}_3$ 线性相关．具体地，可得线性组合关系：
    \begin{equation}
        4(\symbfit{a} - \symbfit{b} + \symbfit{c}) + (-2\symbfit{a} + 3\symbfit{b} - 2\symbfit{c}) - (2\symbfit{a} - \symbfit{b} + 2\symbfit{c}) = \symbfit{0}
    \end{equation}
    即
    \begin{equation}
        \boxed{4\symbfit{v}_1 + \symbfit{v}_2 - \symbfit{v}_3 = \symbfit{0}}
    \end{equation}
\end{myproof}


\begin{problem}{空间向量分解定理}{Space-Vector-Decomposition}
    证明: 三维空间中四个或四个以上的向量一定线性相关.
\end{problem}

\begin{myproof}
    设 $\symbfit{a}, \symbfit{b}, \symbfit{c}, \symbfit{d}$ 为三维空间中的任意四个向量．

    若 $\symbfit{a}, \symbfit{b}, \symbfit{c}$ 中存在共面（或共线）的情况，则该子向量组已线性相关，从而整体必线性相关．

    若 $\symbfit{a}, \symbfit{b}, \symbfit{c}$ 不共面，根据空间向量基本定理，空间中任意向量 $\symbfit{d}$ 均可由这三个不共面向量线性表示．即存在实数 $\lambda_1, \lambda_2, \lambda_3$ 使得
    \begin{equation}
        \symbfit{d} = \lambda_1\symbfit{a} + \lambda_2\symbfit{b} + \lambda_3\symbfit{c}
    \end{equation}
    由向量运算性质可得
    \begin{equation}
        \lambda_1\symbfit{a} + \lambda_2\symbfit{b} + \lambda_3\symbfit{c} + (-1)\symbfit{d} = \symbfit{0}
    \end{equation}
    由于向量 $\symbfit{d}$ 的系数 $-1 \neq 0$，说明存在不全为零的常数使得该向量组的线性组合为零向量．因此，这四个向量线性相关．

    对于四个以上的向量，由于其包含的四个向量子组已经线性相关，根据线性相关性的性质，增加向量后的向量组
    \begin{equation}
        \boxed{\text{一定线性相关}}
    \end{equation}
\end{myproof}


\begin{problem}{向量基底与共面}{Vector-Basis}
    设 $\symbfit{e}_1, \symbfit{e}_2, \symbfit{e}_3$ 为一组基，
    \begin{enumerate}
        \item 证明：$\symbfit{a} = \symbfit{e}_1 + 2\symbfit{e}_2 - \symbfit{e}_3$, $\symbfit{b} = 2\symbfit{e}_1 + \symbfit{e}_2 + \symbfit{e}_3$, $\symbfit{c} = 3\symbfit{e}_1 + 2\symbfit{e}_3$ 为一组基；
        \item 设 $\tilde{\symbfit{c}} = 3\symbfit{e}_1 + x\symbfit{e}_2 + 2\symbfit{e}_3$. 问当 $x$ 取何值时，$\symbfit{a}, \symbfit{b}, \tilde{\symbfit{c}}$ 共面？
    \end{enumerate}
\end{problem}

\begin{solution}
    \ansitem{1}

    由于 $\symbfit{e}_1, \symbfit{e}_2, \symbfit{e}_3$ 为一组基，向量 $\symbfit{a}, \symbfit{b}, \symbfit{c}$ 在该基下的坐标构成的行列式为
    \begin{equation}
        D = \begin{vmatrix} 1 & 2 & 3 \\ 2 & 1 & 0 \\ -1 & 1 & 2 \end{vmatrix}
    \end{equation}
    按第二行展开得
    \begin{equation}
        D = -2 \cdot \begin{vmatrix} 2 & 3 \\ 1 & 2 \end{vmatrix} + 1 \cdot \begin{vmatrix} 1 & 3 \\ -1 & 2 \end{vmatrix} = -2(4-3) + (2+3) = 3
    \end{equation}
    因为 $D = 3 \neq 0$，所以向量组 $\symbfit{a}, \symbfit{b}, \symbfit{c}$ 线性无关．
    \begin{equation}
        \boxed{\text{故 } \symbfit{a}, \symbfit{b}, \symbfit{c} \text{ 构成一组基}}
    \end{equation}

    \ansitem{2}

    若向量 $\symbfit{a}, \symbfit{b}, \tilde{\symbfit{c}}$ 共面，则其对应的坐标行列式必为零
    \begin{equation}
        D_x = \begin{vmatrix} 1 & 2 & 3 \\ 2 & 1 & x \\ -1 & 1 & 2 \end{vmatrix} = 0
    \end{equation}
    按第三行展开计算行列式
    \begin{equation}
        \begin{aligned}
            D_x & = -1 \cdot \begin{vmatrix} 2 & 3 \\ 1 & x \end{vmatrix} - 1 \cdot \begin{vmatrix} 1 & 3 \\ 2 & x \end{vmatrix} + 2 \cdot \begin{vmatrix} 1 & 2 \\ 2 & 1 \end{vmatrix} \\
                & = -(2x - 3) - (x - 6) + 2(1 - 4)                                                                                                                                      \\
                & = -2x + 3 - x + 6 - 6                                                                                                                                                 \\
                & = 3 - 3x
        \end{aligned}
    \end{equation}
    令 $3 - 3x = 0$，解得
    \begin{equation}
        \boxed{x = 1}
    \end{equation}
\end{solution}


\begin{problem}{向量积运算}{Cross-Product-Calculation}
    设向量 $\symbfit{a}, \symbfit{b}$ 的夹角为 $60^\circ$，且 $|\symbfit{a}| = 1, |\symbfit{b}| = 2$. 试求 $|\symbfit{a} \times \symbfit{b}|^2$ 和 $|(\symbfit{a} + \symbfit{b}) \times (\symbfit{a} - \symbfit{b})|$.
\end{problem}

\begin{solution}
    首先计算 $|\symbfit{a} \times \symbfit{b}|$
    \begin{equation}
        |\symbfit{a} \times \symbfit{b}| = |\symbfit{a}| |\symbfit{b}| \sin 60^\circ = 1 \cdot 2 \cdot \frac{\sqrt{3}}{2} = \sqrt{3}
    \end{equation}
    从而有
    \begin{equation}
        \boxed{|\symbfit{a} \times \symbfit{b}|^2 = 3}
    \end{equation}
    利用向量积的分配律与性质 $\symbfit{a} \times \symbfit{a} = \symbfit{0}, \symbfit{b} \times \symbfit{a} = - \symbfit{a} \times \symbfit{b}$，可得
    \begin{equation}
        \begin{aligned}
            (\symbfit{a} + \symbfit{b}) \times (\symbfit{a} - \symbfit{b}) & = \symbfit{a} \times \symbfit{a} - \symbfit{a} \times \symbfit{b} + \symbfit{b} \times \symbfit{a} - \symbfit{b} \times \symbfit{b} \\
                                                                           & = \symbfit{0} - \symbfit{a} \times \symbfit{b} - \symbfit{a} \times \symbfit{b} - \symbfit{0}                                       \\
                                                                           & = -2 (\symbfit{a} \times \symbfit{b})
        \end{aligned}
    \end{equation}
    对其取模
    \begin{equation}
        |(\symbfit{a} + \symbfit{b}) \times (\symbfit{a} - \symbfit{b})| = |-2 (\symbfit{a} \times \symbfit{b})| = 2 |\symbfit{a} \times \symbfit{b}| = 2\sqrt{3}
    \end{equation}
    即
    \begin{equation}
        \boxed{|(\symbfit{a} + \symbfit{b}) \times (\symbfit{a} - \symbfit{b})| = 2\sqrt{3}}
    \end{equation}
\end{solution}


\begin{problem}{雅可比恒等式}{Jacobi-Identity}
    证明不等式：$(\symbfit{a} \times \symbfit{b}) \times \symbfit{c} + (\symbfit{b} \times \symbfit{c}) \times \symbfit{a} + (\symbfit{c} \times \symbfit{a}) \times \symbfit{b} = \symbfit{0}$.
\end{problem}

\begin{myproof}
    根据向量三重积展开公式 $\symbfit{A} \times (\symbfit{B} \times \symbfit{C}) = (\symbfit{A} \cdot \symbfit{C})\symbfit{B} - (\symbfit{A} \cdot \symbfit{B})\symbfit{C}$，以及向量积的反对称性 $(\symbfit{u} \times \symbfit{v}) = - (\symbfit{v} \times \symbfit{u})$，可知
    \begin{equation}
        (\symbfit{a} \times \symbfit{b}) \times \symbfit{c} = - \symbfit{c} \times (\symbfit{a} \times \symbfit{b}) = - [(\symbfit{c} \cdot \symbfit{b})\symbfit{a} - (\symbfit{c} \cdot \symbfit{a})\symbfit{b}] = (\symbfit{c} \cdot \symbfit{a})\symbfit{b} - (\symbfit{c} \cdot \symbfit{b})\symbfit{a}
    \end{equation}
    同理可得剩余两项的展开式
    \begin{equation}
        (\symbfit{b} \times \symbfit{c}) \times \symbfit{a} = (\symbfit{a} \cdot \symbfit{b})\symbfit{c} - (\symbfit{a} \cdot \symbfit{c})\symbfit{b}
    \end{equation}
    \begin{equation}
        (\symbfit{c} \times \symbfit{a}) \times \symbfit{b} = (\symbfit{b} \cdot \symbfit{c})\symbfit{a} - (\symbfit{b} \cdot \symbfit{a})\symbfit{c}
    \end{equation}
    将上述三式相加，并利用数量积的交换律 $\symbfit{u} \cdot \symbfit{v} = \symbfit{v} \cdot \symbfit{u}$
    \begin{equation}
        \begin{aligned}
             & \quad (\symbfit{a} \times \symbfit{b}) \times \symbfit{c} + (\symbfit{b} \times \symbfit{c}) \times \symbfit{a} + (\symbfit{c} \times \symbfit{a}) \times \symbfit{b}                                                                                                               \\
             & = [(\symbfit{c} \cdot \symbfit{a})\symbfit{b} - (\symbfit{c} \cdot \symbfit{b})\symbfit{a}] + [(\symbfit{a} \cdot \symbfit{b})\symbfit{c} - (\symbfit{a} \cdot \symbfit{c})\symbfit{b}] + [(\symbfit{b} \cdot \symbfit{c})\symbfit{a} - (\symbfit{b} \cdot \symbfit{a})\symbfit{c}] \\
             & = (\symbfit{c} \cdot \symbfit{a} - \symbfit{a} \cdot \symbfit{c})\symbfit{b} + (\symbfit{b} \cdot \symbfit{c} - \symbfit{c} \cdot \symbfit{b})\symbfit{a} + (\symbfit{a} \cdot \symbfit{b} - \symbfit{b} \cdot \symbfit{a})\symbfit{c}                                              \\
             & = 0\symbfit{b} + 0\symbfit{a} + 0\symbfit{c}                                                                                                                                                                                                                                        \\
             & = \symbfit{0}
        \end{aligned}
    \end{equation}
    证毕．
\end{myproof}


\begin{problem}{柱面方程}{Cylinder-Equation}
    求准线为 $\begin{cases} y^2 + z^2 = 1, \\ x = 1, \end{cases}$ 母线方向为 $(2, 1, 1)$ 的柱面的一般方程.
\end{problem}

\begin{solution}
    设 $(x, y, z)$ 为柱面上任一点，其母线方向向量为 $\symbfit{s} = (2, 1, 1)$，则过该点的母线方程为
    \begin{equation}
        \frac{X - x}{2} = \frac{Y - y}{1} = \frac{Z - z}{1} = t
    \end{equation}
    母线与准线的交点 $(X, Y, Z)$ 满足
    \begin{equation}
        \begin{cases}
            X = x + 2t = 1 \\
            Y = y + t      \\
            Z = z + t
        \end{cases}
    \end{equation}
    由第一式解得 $t = \frac{1 - x}{2}$，将其代入后两式可得
    \begin{equation}
        Y = y + \frac{1 - x}{2} = \frac{2y - x + 1}{2}, \quad Z = z + \frac{1 - x}{2} = \frac{2z - x + 1}{2}
    \end{equation}
    由于交点在准线 $Y^2 + Z^2 = 1$ 上，故有
    \begin{equation}
        \left(\frac{x - 2y - 1}{2}\right)^2 + \left(\frac{x - 2z - 1}{2}\right)^2 = 1 \implies (x - 2y - 1)^2 + (x - 2z - 1)^2 = 4
    \end{equation}
    展开整理得柱面的一般方程为
    \begin{equation}
        \boxed{x^2 + 2y^2 + 2z^2 - 2xy - 2xz - 2x + 2y + 2z - 1 = 0}
    \end{equation}
\end{solution}


\begin{problem}{锥面方程}{Cone-Equation}
    求准线为 $\begin{cases} y^2 + z^2 = 1, \\ x = 1, \end{cases}$ 顶点坐标为 $(2, 1, 1)$ 的锥面的一般方程.
\end{problem}

\begin{solution}
    设 $(x, y, z)$ 为锥面上除顶点外的任一点，连接顶点 $V(2, 1, 1)$ 与该点的直线（母线）方程为
    \begin{equation}
        \frac{X - 2}{x - 2} = \frac{Y - 1}{y - 1} = \frac{Z - 1}{z - 1} = \lambda
    \end{equation}
    该母线与准线 $X = 1$ 相交，则有
    \begin{equation}
        \frac{1 - 2}{x - 2} = \lambda \implies \lambda = \frac{1}{2 - x}
    \end{equation}
    由此可确定交点坐标为
    \begin{equation}
        Y = \lambda(y - 1) + 1 = \frac{y - x + 1}{2 - x}, \quad Z = \lambda(z - 1) + 1 = \frac{z - x + 1}{2 - x}
    \end{equation}
    将 $Y, Z$ 代入准线方程 $Y^2 + Z^2 = 1$ 得
    \begin{equation}
        \left(\frac{y - x + 1}{2 - x}\right)^2 + \left(\frac{z - x + 1}{2 - x}\right)^2 = 1 \implies (x - y - 1)^2 + (x - z - 1)^2 = (x - 2)^2
    \end{equation}
    展开各项并合并同类项
    \begin{equation}
        (x^2 + y^2 + 1 - 2xy - 2x + 2y) + (x^2 + z^2 + 1 - 2xz - 2x + 2z) = x^2 - 4x + 4
    \end{equation}
    整理得锥面的一般方程为
    \begin{equation}
        \boxed{x^2 + y^2 + z^2 - 2xy - 2xz + 2y + 2z - 2 = 0}
    \end{equation}
\end{solution}


\begin{problem}{旋转面方程}{Surface-Of-Revolution-1}
    求直线 $x - 1 = y = z$ 绕 $x = y = 1$ 旋转所得旋转面的参数方程和一般方程．
\end{problem}

\begin{solution}
    旋转轴为直线 $L: \frac{x-1}{0} = \frac{y-1}{0} = \frac{z}{1}$，其方向向量为 $\symbfit{k} = (0, 0, 1)$．
    设母线上任一点为 $P_0(x_0, y_0, z_0)$，满足：
    \begin{equation}
        x_0 = z_0 + 1, \quad y_0 = z_0.
    \end{equation}
    点 $P(x, y, z)$ 在旋转面上，其到旋转轴 $x=1, y=1$ 的距离等于母线上的点到轴的距离，且 $z$ 坐标保持不变：
    \begin{equation}
        \begin{aligned}
            (x - 1)^2 + (y - 1)^2          & = (x_0 - 1)^2 + (y_0 - 1)^2     \\
            \implies (x - 1)^2 + (y - 1)^2 & = (z_0 + 1 - 1)^2 + (z_0 - 1)^2 \\
                                           & = z_0^2 + (z_0 - 1)^2.
        \end{aligned}
    \end{equation}
    由于 $z = z_0$，代入上式整理得该旋转面的一般方程：
    \begin{equation}
        \boxed{(x - 1)^2 + (y - 1)^2 = 2z^2 - 2z + 1}
    \end{equation}
    令 $z = u, \theta = v$，则旋转面的参数方程可表示为：
    \begin{equation}
        \boxed{\begin{cases}
            x = 1 + \sqrt{2u^2 - 2u + 1} \cos v \\
            y = 1 + \sqrt{2u^2 - 2u + 1} \sin v \\
            z = u
        \end{cases} \quad (u \in \symbb{R}, v \in [0, 2\uppi))}
    \end{equation}
\end{solution}


\begin{problem}{旋转面方程}{Surface-Of-Revolution-2}
    求圆 $\begin{cases} (x - 2)^2 + y^2 = 1, \\ z = 0 \end{cases}$ 绕 $y$ 轴旋转所得旋转面的参数方程和一般方程．
\end{problem}

\begin{solution}
    母线为 $Oxy$ 平面上的圆 $(x - 2)^2 + y^2 = 1$．将其绕 $y$ 轴旋转时，空间中点 $(x, y, z)$ 满足原母线方程中的 $x$ 被替换为 $\pm \sqrt{x^2 + z^2}$．
    代入母线方程得：
    \begin{equation}
        (\sqrt{x^2 + z^2} - 2)^2 + y^2 = 1.
    \end{equation}
    整理得旋转面的一般方程为：
    \begin{equation}
        \boxed{x^2 + y^2 + z^2 - 4\sqrt{x^2 + z^2} + 3 = 0}
    \end{equation}
    为求参数方程，先写出母线圆在 $Oxy$ 平面上的参数表示：
    \begin{equation}
        \begin{cases}
            x = 2 + \cos u \\
            y = \sin u
        \end{cases} \quad (u \in [0, 2\uppi))
    \end{equation}
    设绕 $y$ 轴旋转的角度为 $v$，则旋转面上任意一点的坐标为：
    \begin{equation}
        \boxed{\begin{cases}
            x = (2 + \cos u) \cos v \\
            y = \sin u              \\
            z = (2 + \cos u) \sin v
        \end{cases} \quad (u, v \in [0, 2\uppi))}
    \end{equation}
\end{solution}


\begin{problem}{坐标变换}{Coordinate-Transformation-2}
    将坐标系 $[O; \symbfit{e}_1, \symbfit{e}_2, \symbfit{e}_3]$ 绕 $\symbfit{e}_1$ 逆时针旋转角度 $\alpha$ 后得到新坐标系 $[O; \tilde{\symbfit{e}}_1, \tilde{\symbfit{e}}_2, \tilde{\symbfit{e}}_3]$，再将坐标系 $[O; \tilde{\symbfit{e}}_1, \tilde{\symbfit{e}}_2, \tilde{\symbfit{e}}_3]$ 绕 $\tilde{\symbfit{e}}_2$ 逆时针旋转角度 $\beta$ 后得到的新坐标系为 $[O; \hat{\symbfit{e}}_1, \hat{\symbfit{e}}_2, \hat{\symbfit{e}}_3]$．求 $[O; \symbfit{e}_1, \symbfit{e}_2, \symbfit{e}_3]$ 与 $[O; \hat{\symbfit{e}}_1, \hat{\symbfit{e}}_2, \hat{\symbfit{e}}_3]$ 之间的坐标变换公式．
\end{problem}

\begin{solution}
    设点在 $[O; \symbfit{e}_1, \symbfit{e}_2, \symbfit{e}_3]$ 下的坐标为 $\symbfit{x} = (x_1, x_2, x_3)^{\symrm{T}}$，
    在 $[O; \tilde{\symbfit{e}}_1, \tilde{\symbfit{e}}_2, \tilde{\symbfit{e}}_3]$ 下的坐标为 $\tilde{\symbfit{x}} = (\tilde{x}_1, \tilde{x}_2, \tilde{x}_3)^{\symrm{T}}$，
    在 $[O; \hat{\symbfit{e}}_1, \allowbreak \hat{\symbfit{e}}_2, \allowbreak \hat{\symbfit{e}}_3]$ 下的坐标为 $\hat{\symbfit{x}} = (\hat{x}_1, \hat{x}_2, \hat{x}_3)^{\symrm{T}}$．

    根据绕坐标轴旋转的定义，第一次绕 $\symbfit{e}_1$ 旋转的坐标变换矩阵为
    \begin{equation}
        \symbfit{M}_1 = \begin{bmatrix}
            1 & 0          & 0           \\
            0 & \cos\alpha & -\sin\alpha \\
            0 & \sin\alpha & \cos\alpha
        \end{bmatrix}
    \end{equation}
    则有 $\symbfit{x} = \symbfit{M}_1 \tilde{\symbfit{x}}$．第二次绕新系的 $\tilde{\symbfit{e}}_2$ 旋转的坐标变换矩阵为
    \begin{equation}
        \symbfit{M}_2 = \begin{bmatrix}
            \cos\beta  & 0 & \sin\beta \\
            0          & 1 & 0         \\
            -\sin\beta & 0 & \cos\beta
        \end{bmatrix}
    \end{equation}
    则有 $\tilde{\symbfit{x}} = \symbfit{M}_2 \hat{\symbfit{x}}$．由此可得总的坐标变换公式为
    \begin{equation}
        \begin{aligned}
            \symbfit{x} & = \symbfit{M}_1 \symbfit{M}_2 \hat{\symbfit{x}}                \\
                        & = \begin{bmatrix}
                                1 & 0          & 0           \\
                                0 & \cos\alpha & -\sin\alpha \\
                                0 & \sin\alpha & \cos\alpha
                            \end{bmatrix} \begin{bmatrix}
                                              \cos\beta  & 0 & \sin\beta \\
                                              0          & 1 & 0         \\
                                              -\sin\beta & 0 & \cos\beta
                                          \end{bmatrix} \hat{\symbfit{x}}                \\
                        & = \begin{bmatrix}
                                \cos\beta            & 0          & \sin\beta            \\
                                \sin\alpha\sin\beta  & \cos\alpha & -\sin\alpha\cos\beta \\
                                -\cos\alpha\sin\beta & \sin\alpha & \cos\alpha\cos\beta
                            \end{bmatrix} \hat{\symbfit{x}}
        \end{aligned}
    \end{equation}
    即坐标变换公式为
    \begin{equation}
        \boxed{\begin{bmatrix} x_1 \\ x_2 \\ x_3 \end{bmatrix} = \begin{bmatrix} \cos\beta & 0 & \sin\beta \\ \sin\alpha\sin\beta & \cos\alpha & -\sin\alpha\cos\beta \\ -\cos\alpha\sin\beta & \sin\alpha & \cos\alpha\cos\beta \end{bmatrix} \begin{bmatrix} \hat{x}_1 \\ \hat{x}_2 \\ \hat{x}_3 \end{bmatrix}}
    \end{equation}
\end{solution}


\begin{problem}{二次曲面方程}{Quadric-Surface-Equation}
    已知椭球面的三个半轴长分别为 $a,b,c$，三条对称轴方程分别为
    \begin{equation}
        3 - x = \frac{y}{2} = \frac{z}{2}, \quad \frac{x}{2} = 3 - y = \frac{z}{2}, \quad \frac{x}{2} = \frac{y}{2} = 3 - z,
    \end{equation}
    求椭球面的一般方程．
\end{problem}

\begin{solution}
    首先，将给定的三条对称轴方程化为标准形式：
    \begin{equation}
        \begin{aligned}
            L_1 & \colon \frac{x - 3}{-1} = \frac{y}{2} = \frac{z}{2}, \\
            L_2 & \colon \frac{x}{2} = \frac{y - 3}{-1} = \frac{z}{2}, \\
            L_3 & \colon \frac{x}{2} = \frac{y}{2} = \frac{z - 3}{-1}.
        \end{aligned}
    \end{equation}

    由此可得，这三条对称轴的方向向量依次为
    \begin{equation}
        \symbfit{v}_1 = (-1, 2, 2), \quad \symbfit{v}_2 = (2, -1, 2), \quad \symbfit{v}_3 = (2, 2, -1).
    \end{equation}

    易验证这三个方向向量两两正交，即 $\symbfit{v}_1 \cdot \symbfit{v}_2 = 0$, $\symbfit{v}_2 \cdot \symbfit{v}_3 = 0$, $\symbfit{v}_3 \cdot \symbfit{v}_1 = 0$．且它们的模长均为
    \begin{equation}
        \|\symbfit{v}_1\| = \|\symbfit{v}_2\| = \|\symbfit{v}_3\| = \sqrt{(-1)^2 + 2^2 + 2^2} = 3.
    \end{equation}

    因此，对应的单位方向向量为
    \begin{equation}
        \symbfit{e}_1 = \frac{1}{3}(-1, 2, 2), \quad \symbfit{e}_2 = \frac{1}{3}(2, -1, 2), \quad \symbfit{e}_3 = \frac{1}{3}(2, 2, -1).
    \end{equation}

    其次，求椭球面的对称中心．三条对称轴的交点即为椭球面的中心 $M_0(x_0, y_0, z_0)$．联立 $L_1$ 与 $L_2$ 的方程：
    \begin{equation}
        3 - x = \frac{y}{2} = \frac{z}{2} \quad \text{与} \quad \frac{x}{2} = 3 - y = \frac{z}{2}.
    \end{equation}

    由 $\frac{z}{2} = \frac{z}{2}$ 可设公共参数为 $t$，则 $z = 2t$．代入得 $x = 3 - t$ 且 $y = 2t$，同时从 $L_2$ 中有 $\frac{3 - t}{2} = 3 - 2t \implies t = 1$．
    将 $t = 1$ 代回，解得交点坐标为 $(2, 2, 2)$．显然该点也满足 $L_3$ 的方程，故椭球面的中心为 $M_0(2, 2, 2)$．

    设空间任一点 $P(x, y, z)$ 在以 $M_0$ 为原点，以 $\symbfit{e}_1, \symbfit{e}_2, \symbfit{e}_3$ 为基向量的新直角坐标系中的坐标为 $(X,\allowbreak Y,\allowbreak Z)$．记向量 $\overrightarrow{M_0 P} = (x - 2, y - 2, z - 2)$，则由坐标投影关系可得：
    \begin{equation}
        \begin{aligned}
            X & = \overrightarrow{M_0 P} \cdot \symbfit{e}_1 = \frac{-(x - 2) + 2(y - 2) + 2(z - 2)}{3} = \frac{-x + 2y + 2z - 6}{3}, \\
            Y & = \overrightarrow{M_0 P} \cdot \symbfit{e}_2 = \frac{2(x - 2) - (y - 2) + 2(z - 2)}{3} = \frac{2x - y + 2z - 6}{3},   \\
            Z & = \overrightarrow{M_0 P} \cdot \symbfit{e}_3 = \frac{2(x - 2) + 2(y - 2) - (z - 2)}{3} = \frac{2x + 2y - z - 6}{3}.
        \end{aligned}
    \end{equation}

    在新坐标系下，由于三个半轴长分别为 $a, b, c$，椭球面的标准方程为
    \begin{equation}
        \frac{X^2}{a^2} + \frac{Y^2}{b^2} + \frac{Z^2}{c^2} = 1.
    \end{equation}

    将上述 $X, Y, Z$ 的表达式代入，即可得到原坐标系下的椭球面一般方程：
    \begin{equation}
        \boxed{\frac{(-x + 2y + 2z - 6)^2}{9a^2} + \frac{(2x - y + 2z - 6)^2}{9b^2} + \frac{(2x + 2y - z - 6)^2}{9c^2} = 1}
    \end{equation}
\end{solution}


\begin{problem}{球面坐标系两点距离}{Spherical-Distance}
    设在球面坐标系中有两点 $(r_1, \theta_1, \varphi_1)$ 和 $(r_2, \theta_2, \varphi_2)$， $d$ 是这两点之间的直线距离. 证明：
    \begin{equation}
        d = \sqrt{(r_1 - r_2)^2 + 2r_1 r_2 [1 - \cos(\varphi_1 - \varphi_2) \sin\theta_1 \sin\theta_2 - \cos\theta_1 \cos\theta_2]}.
    \end{equation}
\end{problem}

\begin{myproof}
    设这两点对应的位置向量分别为 $\symbfit{r}_1$ 和 $\symbfit{r}_2$．在直角坐标系下，根据球面坐标与直角坐标的转换关系，有
    \begin{equation}
        \begin{aligned}
            \symbfit{r}_1 & = (r_1 \sin\theta_1 \cos\varphi_1, r_1 \sin\theta_1 \sin\varphi_1, r_1 \cos\theta_1) \\
            \symbfit{r}_2 & = (r_2 \sin\theta_2 \cos\varphi_2, r_2 \sin\theta_2 \sin\varphi_2, r_2 \cos\theta_2)
        \end{aligned}
    \end{equation}

    两点之间的直线距离的平方等于其位置向量差的模的平方，即
    \begin{equation}
        d^2 = |\symbfit{r}_1 - \symbfit{r}_2|^2 = |\symbfit{r}_1|^2 + |\symbfit{r}_2|^2 - 2 \symbfit{r}_1 \cdot \symbfit{r}_2
    \end{equation}

    易知 $|\symbfit{r}_1|^2 = r_1^2$，$|\symbfit{r}_2|^2 = r_2^2$．计算两向量的内积可得
    \begin{equation}
        \begin{aligned}
            \symbfit{r}_1 \cdot \symbfit{r}_2 & = (r_1 \sin\theta_1 \cos\varphi_1)(r_2 \sin\theta_2 \cos\varphi_2) + (r_1 \sin\theta_1 \sin\varphi_1)(r_2 \sin\theta_2 \sin\varphi_2) + (r_1 \cos\theta_1)(r_2 \cos\theta_2) \\
                                              & = r_1 r_2 [\sin\theta_1 \sin\theta_2 (\cos\varphi_1 \cos\varphi_2 + \sin\varphi_1 \sin\varphi_2) + \cos\theta_1 \cos\theta_2]                                                \\
                                              & = r_1 r_2 [\cos(\varphi_1 - \varphi_2) \sin\theta_1 \sin\theta_2 + \cos\theta_1 \cos\theta_2]
        \end{aligned}
    \end{equation}

    将上述结果代入距离平方公式中，得到
    \begin{equation}
        d^2 = r_1^2 + r_2^2 - 2r_1 r_2 [\cos(\varphi_1 - \varphi_2) \sin\theta_1 \sin\theta_2 + \cos\theta_1 \cos\theta_2]
    \end{equation}

    利用代数恒等式 $(r_1 - r_2)^2 = r_1^2 + r_2^2 - 2r_1 r_2$ 将上式进行配凑变形：
    \begin{equation}
        \begin{aligned}
            d^2 & = (r_1^2 - 2r_1 r_2 + r_2^2) + 2r_1 r_2 - 2r_1 r_2 [\cos(\varphi_1 - \varphi_2) \sin\theta_1 \sin\theta_2 + \cos\theta_1 \cos\theta_2] \\
                & = (r_1 - r_2)^2 + 2r_1 r_2 [1 - \cos(\varphi_1 - \varphi_2) \sin\theta_1 \sin\theta_2 - \cos\theta_1 \cos\theta_2]
        \end{aligned}
    \end{equation}

    由于 $d \ge 0$，两边同时开平方即得所求证的距离公式
    \begin{equation}
        \boxed{d = \sqrt{(r_1 - r_2)^2 + 2r_1 r_2 [1 - \cos(\varphi_1 - \varphi_2) \sin\theta_1 \sin\theta_2 - \cos\theta_1 \cos\theta_2]}}
    \end{equation}
\end{myproof}

\begin{problem}{球面距离}{Spherical-Surface-Distance}
    设在球面 $r=a$ 上有两个点 $(a, \theta_1, \varphi_1)$ 和 $(a, \theta_2, \varphi_2)$. 证明：这两个点之间的球面距离是 $a\gamma$ $(0 \le \gamma \le \pi)$，其中 $\gamma$ 满足
    \begin{equation}
        \cos\gamma = \cos(\varphi_1 - \varphi_2)\sin\theta_1\sin\theta_2 + \cos\theta_1\cos\theta_2.
    \end{equation}
\end{problem}

\begin{myproof}
    设球心为原点 $O$，两点对应的位置向量分别为 $\symbfit{r}_1$ 和 $\symbfit{r}_2$．由于这两点均位于半径为 $a$ 的球面上，它们在直角坐标系中的位置向量可表示为
    \begin{equation}
        \begin{aligned}
            \symbfit{r}_1 & = a (\sin\theta_1 \cos\varphi_1, \sin\theta_1 \sin\varphi_1, \cos\theta_1) \\
            \symbfit{r}_2 & = a (\sin\theta_2 \cos\varphi_2, \sin\theta_2 \sin\varphi_2, \cos\theta_2)
        \end{aligned}
    \end{equation}

    球面距离被定义为经过这两点的大圆劣弧的长度．设向量 $\symbfit{r}_1$ 和 $\symbfit{r}_2$ 之间的夹角为 $\gamma$ $(0 \le \gamma \le \pi)$，则由弧长公式可知，该球面距离即为 $a\gamma$．

    根据向量内积的几何定义，有
    \begin{equation}
        \symbfit{r}_1 \cdot \symbfit{r}_2 = |\symbfit{r}_1| |\symbfit{r}_2| \cos\gamma = a^2 \cos\gamma
    \end{equation}

    另一方面，根据向量内积的坐标运算，可得
    \begin{equation}
        \begin{aligned}
            \symbfit{r}_1 \cdot \symbfit{r}_2 & = a^2 (\sin\theta_1 \cos\varphi_1 \sin\theta_2 \cos\varphi_2 + \sin\theta_1 \sin\varphi_1 \sin\theta_2 \sin\varphi_2 + \cos\theta_1 \cos\theta_2) \\
                                              & = a^2 [\sin\theta_1 \sin\theta_2 (\cos\varphi_1 \cos\varphi_2 + \sin\varphi_1 \sin\varphi_2) + \cos\theta_1 \cos\theta_2]                         \\
                                              & = a^2 [\cos(\varphi_1 - \varphi_2) \sin\theta_1 \sin\theta_2 + \cos\theta_1 \cos\theta_2]
        \end{aligned}
    \end{equation}

    将上述两种内积表达方式的等号相连，由于 $a \neq 0$，等式两边同除以 $a^2$ 即可得出夹角 $\gamma$ 所满足的等式关系：
    \begin{equation}
        \boxed{\cos\gamma = \cos(\varphi_1 - \varphi_2) \sin\theta_1 \sin\theta_2 + \cos\theta_1 \cos\theta_2}
    \end{equation}
\end{myproof}

\endinput
